% !TEX program = xelatex
% ============================================================
%  SY02-46  Drinfeld–Langlands 对应：Frenkel、Gaitsgory 与 Vilonen 的工作
%  Gérard Laumon，Séminaire BOURBAKI，54e année (2001–2002)，n° 906，Juin 2002
%  编译：xelatex 三遍
% ============================================================
\documentclass[11pt,a4paper]{article}
% ============================================================
%  SY02-46  Laumon, Travaux de Frenkel, Gaitsgory et Vilonen
%           sur la correspondance de Drinfeld-Langlands
%  Séminaire Bourbaki n° 906 (Juin 2002)  —— 中译本 preamble
% ============================================================
\usepackage[UTF8,fontset=windows,heading=false]{ctex}
\usepackage{amsmath,amssymb,amsthm,mathtools,bm}
\usepackage[margin=2.6cm]{geometry}
\usepackage{xcolor}
\usepackage[hidelinks]{hyperref}
\usepackage{tikz-cd}
\usepackage{array}

% ---------- 版式 ----------
\linespread{1.18}
\setlength{\parskip}{0.25em}
\setlength{\marginparwidth}{2.1cm}
\renewcommand{\proofname}{证明}
\renewcommand{\qedsymbol}{$\square$}

% ---------- 编号环境 ----------
% 原文体例：类型名在前（THÉORÈME 2.2. — …），同一节内定理、引理、定义、公式共用
% 一个序号；故让全部环境共享 [section] 计数器，按原文出现顺序排版即可还原编号。
\theoremstyle{definition}
\newtheorem{theorem}{定理}[section]
\newtheorem{definition}[theorem]{定义}
\newtheorem{lemma}[theorem]{引理}
\newtheorem{proposition}[theorem]{命题}
\newtheorem{corollary}[theorem]{推论}
\newtheorem{remark}[theorem]{注}

% §0 的条目在原文中编号为 (0.1)…(0.4)，且与公式共用序号；
% 这里用不编号环境 + 可选参数手动给出括号号，避免与第 1 节的自动编号冲突。
\newtheorem*{definitionZ}{定义}
\newtheorem*{lemmaZ}{引理}

% ---------- 原文页码锚（页眉 906–NN，与 PDF 页一致）----------
\newcommand{\srcpage}[1]{\marginpar{\raggedright\footnotesize\ttfamily\color{gray}906--#1}}

% ---------- 译注：\yizhu{内容}{等级（出处）} ----------
% 等级只许三类：术语背景 / 原文省略 / 原文笔误；出处只许 A·B·I 或「上下文推断」
\newcommand{\yizhu}[2]{\footnote{\textsf{【译注·#2】}#1}}

% ---------- 常用对象 ----------
\newcommand{\Fib}{\operatorname{Fib}}
\newcommand{\Coh}{\operatorname{Coh}}
\newcommand{\Mod}{\operatorname{Mod}}
\newcommand{\Perv}{\operatorname{Perv}}
\newcommand{\Pic}{\operatorname{Pic}}
\newcommand{\Jac}{\operatorname{Jac}}
\newcommand{\Four}{\operatorname{Four}}
\newcommand{\Rad}{\operatorname{Rad}}
\newcommand{\Hck}{\operatorname{Hecke}}
\newcommand{\Aut}{\operatorname{Aut}}
\newcommand{\Sym}{\operatorname{Sym}}
\newcommand{\Ext}{\operatorname{Ext}}
\newcommand{\Hom}{\operatorname{Hom}}
\newcommand{\Spec}{\operatorname{Spec}}
\newcommand{\Det}{\operatorname{det}}
\newcommand{\id}{\mathrm{id}}
\newcommand{\pt}{\mathrm{pt}}
\newcommand{\rss}{\mathrm{rss}}
\newcommand{\Bun}{\operatorname{Bun}}
\newcommand{\Rep}{\operatorname{Rep}}
\newcommand{\LocSys}{\operatorname{LocSys}}
\newcommand{\tr}{\operatorname{tr}}
\newcommand{\op}{\mathrm{op}}
\newcommand{\der}{\mathrm{der}}

% 上同调与导出范畴
\newcommand{\Db}{D^{b}}
\newcommand{\Dbc}{D^{b}_{c}}

% 花体
\newcommand{\cA}{\mathcal{A}}
\newcommand{\cB}{\mathcal{B}}
\newcommand{\cC}{\mathcal{C}}
\newcommand{\cD}{\mathcal{D}}
\newcommand{\cE}{\mathcal{E}}
\newcommand{\cF}{\mathcal{F}}
\newcommand{\cG}{\mathcal{G}}
\newcommand{\cH}{\mathcal{H}}
\newcommand{\cI}{\mathcal{I}}
\newcommand{\cJ}{\mathcal{J}}
\newcommand{\cK}{\mathcal{K}}
\newcommand{\cL}{\mathcal{L}}
\newcommand{\cM}{\mathcal{M}}
\newcommand{\cN}{\mathcal{N}}
\newcommand{\cO}{\mathcal{O}}
\newcommand{\cP}{\mathcal{P}}
\newcommand{\cQ}{\mathcal{Q}}
\newcommand{\cR}{\mathcal{R}}
\newcommand{\cS}{\mathcal{S}}
\newcommand{\cT}{\mathcal{T}}
\newcommand{\cU}{\mathcal{U}}
\newcommand{\cV}{\mathcal{V}}
\newcommand{\cW}{\mathcal{W}}
\newcommand{\cX}{\mathcal{X}}
\newcommand{\cY}{\mathcal{Y}}
\newcommand{\cZ}{\mathcal{Z}}

% 数系与域
\newcommand{\Z}{\mathbb{Z}}
\newcommand{\Q}{\mathbb{Q}}
\newcommand{\R}{\mathbb{R}}
\newcommand{\C}{\mathbb{C}}
\newcommand{\F}{\mathbb{F}}
\newcommand{\A}{\mathbb{A}}
\newcommand{\Gm}{\mathbb{G}_m}
\newcommand{\Ql}{\overline{\mathbb{Q}}_{\ell}}

% 常见缩写
\newcommand{\OmX}{\Omega^{1}_{X}}
\newcommand{\Dmod}{D\text{-}\mathrm{mod}}
\newcommand{\dR}{\mathrm{dR}}
\newcommand{\et}{\mathrm{et}}
\newcommand{\Betti}{\mathrm{Betti}}
\newcommand{\lisse}{\mathrm{lisse}}
\newcommand{\irred}{\mathrm{irred}}
\newcommand{\Pv}{\mathbb{P}}
\newcommand{\GG}{\mathbb{G}}

% ---------- 排版小工具 ----------
% 参考文献条目
\newcommand{\refitem}[2]{\par\noindent\hangindent=2em\hangafter=1 #1\ #2\par}
% 行内长串不断行
\newcommand{\dsp}{\displaystyle}

% 放宽断行，消除个别 10pt 级别的 Overfull
\emergencystretch=1.5em


\title{Drinfeld--Langlands 对应：\\[0.35em]
Frenkel、Gaitsgory 与 Vilonen 的工作}
\author{Gérard Laumon\\[0.3em]
\small Université Paris-Sud 与 CNRS，UMR 8628\\[0.2em]
\small Séminaire BOURBAKI，第 54 年（2001--2002），第 906 期，2002 年 6 月}
\date{}

\begin{document}
\maketitle

\noindent\small\raggedright
\textbf{底本}：Gérard Laumon, \emph{Travaux de Frenkel, Gaitsgory et Vilonen sur la
correspondance de Drinfeld--Langlands},\allowbreak{} Séminaire BOURBAKI,\allowbreak{}
54e année (2001--2002),\allowbreak{} n° 906,\allowbreak{} Juin 2002,\allowbreak{}
19 页；\allowbreak{} arXiv:math/0207078v1（2002 年 7 月 9 日）。
dvips 排版，页眉编号 \texttt{906--NN} 与 PDF 页码一致，译文边注中的页码即指此页。\par

\medskip
\noindent\small
\textbf{译注约定}：原文含糊、省略或笔误之处，以脚注标出，脚注以「【译注·类别】」起首，
类别只取\emph{术语背景}、\emph{原文省略}、\emph{原文笔误}三种，并注明依据等级
（A：权威文献；B：nLab/Stacks/ProofWiki 一类百科；I：译者按上下文推断）。
译文不补写原文所无的论证。

\medskip
\noindent\small
\textbf{编号}：原文的定理、引理、定义、命题、推论、注与编号公式在同一节内共用序号。
译文沿用原编号（如「定理 2.2」「引理 (0.3)」），第 0 节的条目写作 $(0.n)$，第 $i$ 节
的条目写作 $i.n$，未作改动。

\normalsize
\tableofcontents
\newpage

% s00-intro.tex —— 印刷页 1（引言）
% 引言在原书中无节号；此处用不编号节，避免占用第 1 节号而使 §1 几何类域论
% 及其后各节的定理编号整体后移一位。
\section*{引言}
\addcontentsline{toc}{section}{引言}
\srcpage{1}

1967 年，R. Langlands 提出了 E. Artin 与 J. Tate 的阿贝尔类域论的一个宏大扩展。更确切地说，他猜想：整体域 $F$ 上约化群 $G$ 的自守表示，与 $F$ 的取值于 $G$ 的对偶代数群 ${}^{L}G$ 的 Galois 表示之间，存在自然的对应。${}^{L}G$ 的单位分支 $\widehat{G}$ 是复约化群，其根是 $G$ 的余根，反之亦然。

若 $G$ 是 $F$ 上的线性群 $\mathrm{GL}(n)$，则 $\widehat{G}$ 正是 $\mathrm{GL}(n,\C)$。当 $F$ 是函数域，即 $\F_p(t)$ 的有限扩张时，整体 Langlands 对应已由 V. Drinfeld [Dr 1]（$n=2$）与 L. Lafforgue [La]（$n$ 任意）证明。数域上的整体 Langlands 对应仍是数学中尚未解决的大问题之一。

Drinfeld–Langlands 对应，又称几何 Langlands 对应，是 Langlands 对应在以下情形的猜想类比：$k(t)$ 的有限扩张 $F$ 上的分裂约化群 $G$，其中 $k$ 是任意域。若 $X$ 是 $k$ 上光滑拟射影的代数曲线，函数域为 $F$，则该对应使 $X$ 上 $G$-丛的模空间与 $X$ 上 $\widehat{G}$-局部系统的模空间互为对偶。

当 $G=\mathrm{GL}(1)$ 时，Drinfeld–Langlands 对应正是 M. Rosenlicht 与 S. Lang 的几何类域论，由 J.-P. Serre 在 [Se] 中阐述。$G=\mathrm{GL}(2)$ 的情形由 V. Drinfeld 在 [Dr 2] 与 [Dr 3] 两文中处理，这两篇文章正是该理论的源头。

许多工作研究了 Drinfeld–Langlands 对应的各个方面，特别是 A. Beilinson 与 V. Drinfeld（[B-D]）、A. Braverman 与 D. Gaitsgory（[B-G]）、以及 S. Lysenko（[Ly 1] 与 [Ly 2]）的工作。在本文中，我只讨论 E. Frenkel、D. Gaitsgory 与 K. Vilonen 的近期工作，即 $X$ 射影（处处非分歧对应）且 $G=\mathrm{GL}(n)$ 的情形。这些工作推广了 V. Drinfeld 在 [Dr 2] 中的结果。

我感谢 S. Lysenko 在本文准备过程中给予的帮助。

% s01-prelim.tex —— 印刷页 2–4（§0 预备）
\section*{0.\ 预备}
\addcontentsline{toc}{section}{0.\ 预备}
\srcpage{2}

本文讨论处处非分歧的情形。此时人们希望对每个三元组建立 Drinfeld–Langlands 对应。三元组由底域 $k$、系数域 $C$，以及 $k$ 上有限型概形范畴的 $C$-系数上同调理论组成。

三个主要的三元组（底域、系数域、上同调理论）是：
\begin{itemize}
\item[(Betti)] $k=\C$，$C$ 是特征 $0$ 的代数闭域\yizhu{原文系数域用斜体 $C$ 记，而 Betti 情形的底域取 $k=\C$（复数域）；两者字母相同而含义不同，勿混。}{术语背景，I（原文 p2；上下文推断）}，以及经典拓扑下 $C$-向量空间的可构造层论，
\item[(De Rham)] $k$ 的特征为 $0$，$C=k$，以及 Zariski 拓扑下 holonomic $D$-模的理论，
\item[($\ell$-进)] $k$ 任意，$C=\Ql$，其中素数 $\ell$ 在 $k$ 中可逆，以及 étale 拓扑下 $\ell$-进层的理论。
\end{itemize}

Drinfeld–Langlands 对应的首要数据，是底域 $k$ 上的一条代数曲线。对此曲线，人们附上一些模空间，它们一般是 $k$ 上的代数叠。因此我将自由使用代数叠的语言（参 [L-M]）。

上述上同调理论对代数叠范畴尚未得到系统的建立，文献也零散。对 Frenkel、Gaitsgory 与 Vilonen 在其文章中所用的 $\ell$-进三元组，情况相当令人满意。事实上，他们为获得几何 Fourier 变换还假设 $k$ 的特征为 $p>0$，但这一限制并不必要。Beilinson 与 Drinfeld 的工作 [B-D] 使用 De Rham 三元组，其中可以找到代数叠上微分算子环的定义。Betti 三元组原则上最基本，Bernstein 与 Luntz [B-L] 研究过它，但解析叠的理论尚待写出。

本文仍然使用最后这个 Betti 三元组。因此底域为 $k=\C$。设 $\cS$ 是有限型代数叠，或更一般地是 $\C$ 上局部有限型的代数叠。对与 $\cS$ 相伴的解析叠上具有有界可构造上同调的 $C$-向量空间层复形，其导出范畴我简记为 $\Dbc(\cS)$。当 $\cS$ 表示为概形 $S$ 按代数群 $G$ 的作用所得的商时，该解析叠正是 $S(\C)^{\mathrm{an}}$ 按 $G(\C)^{\mathrm{an}}$ 的作用所得的商叠\yizhu{原文用同一个字母 $S$ 既记代数叠又记用来呈现它的概形，此处按上下文区分。}{术语背景，I（原文 p2；上下文推断）}。此时 $\Dbc(\cS)$ 就是 Bernstein 与 Luntz 在 [B-L] 中引入的导出范畴 $D^{b}_{G(\C)^{\mathrm{an}}}(S(\C)^{\mathrm{an}})$。

对本文考虑的所有代数叠 $\cS$，导出范畴 $\Dbc(\cS)$ 带有张量积与对偶函子 $D:\Dbc(\cS)^{\mathrm{opp}}\to\Dbc(\cS)$。对这类叠之间所考虑的可表态射 $f:\cS\to\cT$，\srcpage{3}有前推函子 $f_*,f_!:\Dbc(\cS)\to\Dbc(\cT)$ 与拉回函子 $f^*,f^!:\Dbc(\cT)\to\Dbc(\cS)$，它们满足 Grothendieck 六种运算的形式体系。范畴 $\Dbc(\cS)$ 带有中间反常 $t$-结构，其心是 $\cS$ 上反常层的范畴 $\Perv(\cS)$（参 [B-B-D]）。回忆：$\Dbc(\cS)$ 的这个满子范畴是阿贝尔的、Noether 的、Artin 的，并且关于对偶 $D$ 自对偶。

设 $\pi:\cV\to\cS$ 是常秩 $r$ 的向量丛。可考虑 $\cV$ 按乘法群 $\Gm$ 的数乘作用所得的 $\cS$-商叠 $\overline{\cV}$。若 $\cV^\circ$ 是 $\cV$ 中零截面的补开集，则包含 $\cV^\circ\subset\cV$ 诱导开浸入
\[ j:\Pv(\cV)=[\cV^\circ/\Gm]=\overline{\cV}{}^\circ\hookrightarrow\overline{\cV}, \]
其中 $\Pv(\cV)\to\cS$ 是 $\cV\to\cS$ 的直线所成的射影丛。该开浸入的余闭子叠，是 $\cV$ 的零截面按 $\Gm$ 的平凡作用所得的商，即分类叠 $B(\Gm/\cS)$。

$\cV$ 上 $\Gm$-等变的 $C$-向量空间可构造复形，按定义就是 $\Dbc(\overline{\cV})$ 的对象。

对这些 $\Gm$-等变复形，存在几何 Fourier 变换，称为与向量丛 $\cV/\cS$ 相伴的齐次 Fourier 变换（参 [Lau 3]）。这个 Fourier 变换是导出范畴的等价
\[ \Four_{\overline{\cV}/\cS}:\Dbc(\overline{\cV})\to\Dbc(\overline{\cV}^\vee), \]
其逆为 $\Four_{\overline{\cV}^\vee/\cS}$。这里 $\overline{\cV}^\vee$ 是对偶向量丛 $\pi^\vee:\cV^\vee\to\cS$ 按乘法群 $\Gm$ 的数乘作用所得的 $\cS$-商叠。它与对偶交换，并且是 $t$-正合的。可以如下定义它。

在商叠 $\cA=[\A^1/\Gm]$ 上，即在仿射直线按乘法群的数乘作用所得的商上，有复形 $\Psi=\beta_*C\in\operatorname{ob}\Dbc(\cA)$。这里 $\beta:\Spec(\C)=[\Gm/\Gm]\hookrightarrow\cA$ 是由包含 $\Gm\subset\A^1$ 诱导的开浸入。此复形的非零上同调层为
\[ \cH^0(\Psi)=C \]
以及
\[ \cH^1(\Psi)=\alpha_*C, \]
其中 $\alpha:B(\Gm)\hookrightarrow\cA$ 是与 $\beta$ 互补的闭浸入，由原点到 $\A^1$ 的包含诱导。实际上 $\Psi$ 不是不可约反常层，它是 $\alpha_*C[-1]$ 通过点状常数反常层 $C$ 的扩张。

向量丛 $\cV$ 与其对偶丛之间的自然配对态射，过商后成为代数叠的态射
\[ \mu:\overline{\cV}^\vee\times_\cS\overline{\cV}\to\cA. \]
\srcpage{4}%
又若记 $p:\overline{\cV}^\vee\times_\cS\overline{\cV}\to\overline{\cV}$ 与 $p^\vee:\overline{\cV}^\vee\times_\cS\overline{\cV}\to\overline{\cV}^\vee$ 为两个典范投影，则按定义
\[ \tag*{(0.1)} \Four_{\overline{\cV}/\cS}(K)=(p^\vee)_!\bigl(p^*K\otimes\mu^*\Psi\bigr)[r-1],\qquad \forall K\in\Dbc(\overline{\cV}). \]

态射 $p$ 与 $p^\vee$ 不是可表态射，且 $(p^\vee)_!$ 不保持 $\Dbc$：它把 $\Dbc(\overline{\cV}^\vee\times_\cS\overline{\cV})$ 送到 $D^-_c(\overline{\cV}^\vee)$ 中。然而可以验证，$\Four_{\overline{\cV}/\cS}$ 本身保持 $\Dbc$。

\begin{definitionZ}[0.2]
设 $L$ 是 $\Pv(\cV)$ 上的反常层。若遗忘支撑的态射 $j_!L\to j_*L$ 是同构，则称 $L$ 是 proper 的\yizhu{此处的 proper 是 (0.2) 给出的反常层术语，与态射的「紧合」（propre）无关，故保留原文不译。}{术语背景，A（原文 (0.2)；参 BBD，Astérisque 100）}。此时 $j_!L\cong j_*L$ 实际上是 $\overline{\cV}$ 上的反常层，它正是 $L$ 的中间延拓 $j_{!*}L$。
\end{definitionZ}

\begin{lemmaZ}[0.3]
设 $L$ 是 $\Pv(\cV)$ 上的反常层。则 $L$ 是 proper 的，当且仅当 $(\bar\pi^\circ)_!L=0$，其中 $\bar\pi^\circ:\Pv(\cV)\to\cS$ 是典范投影。\hfill$\square$
\end{lemmaZ}

与射影丛 $\Pv(\cV)\to\cS$ 相伴的几何 Radon 变换，是函子
\[ \Rad_{\Pv(\cV)/\cS}:\Dbc(\Pv(\cV))\to\Dbc(\Pv(\cV^\vee)), \]
它由
\[ \Rad_{\Pv(\cV)/\cS}(L)=(q^\vee)_!q^*L[r-1] \]
定义。这里 $q:\cH\to\Pv(\cV)$ 与 $q^\vee:\cH\to\Pv(\cV^\vee)$ 是关联叠
\[ \cH\subset\Pv(\cV^\vee)\times_\cS\Pv(\cV) \]
的两个典范投影。该关联叠由 $\cV\to\cS$ 的一条直线与一个超平面构成的对组成，其中直线含于该超平面。关于这个变换的详细研究，参看 Brylinski 的专著 [Br]。

\begin{lemmaZ}[0.4]
设 $L$ 是 $\Pv(\cV)$ 上的反常层，并假设它是 proper 的。则 $\Four_{\overline{\cV}/\cS}(j_!L)$ 也是 proper 的。它在开集 $\Pv(\cV^\vee)=\overline{\cV}^{\vee\circ}$ 上的限制等于 $\Rad_{\Pv(\cV)/\cS}(L)$；该开集是 $\overline{\cV}^\vee$ 中零截面的补集。\hfill$\square$
\end{lemmaZ}

\section{几何类域论}\srcpage{4}

设 $X$ 是连通、射影、光滑的代数曲线（紧连通 Riemann 面），亏格为 $g$。取定基点 $x_0$、自由 $\Z$-模
\[
  H_1(X,\Z)=\pi_1^{\mathrm{ab}}(X):=\pi_1(X,x_0)/[\pi_1(X,x_0),\pi_1(X,x_0)]
\]
的一组基 $(\delta_1,\dots,\delta_{2g})$，以及复向量空间 $H^{0}(X,\OmX)$ 的一组基 $(\omega_1,\dots,\omega_g)$。周期
\[
  \Pi_i=\left(\int_{\delta_i}\omega_1,\dots,\int_{\delta_i}\omega_g\right)\in\C^{g},
  \qquad i=1,\dots,2g,
\]
\srcpage{5}在 $\Z$ 上线性无关，并生成一个格 $\Lambda\subset\C^{g}$。$X$ 的雅可比簇是 $g$ 维复环面
\[
  J(X)=\C^{g}/\Lambda.
\]
有一个解析映射
\[
  \varphi:X\to J(X),\quad
  x\mapsto\left(\int_{x_0}^{x}\omega_1,\dots,\int_{x_0}^{x}\omega_g\right)+\Lambda,
\]
它把 $x_0$ 送到 $J(X)$ 的单位元 $0$。

设 $\Pic(X)$ 是参数化 $X$ 上线丛同构类的 Picard 概形。线丛的张量积给出群结构，使它成为群概形，并落在正合列
\[
  1\to\Pic^{0}(X)\to\Pic(X)\xrightarrow{\deg}\Z\to 0
\]
中；这里 $\deg(\cL)$ 是线丛 $\cL$ 的次数，$\Pic(X)$ 的单位分支 $\Pic^{0}(X)$ 是阿贝尔簇。

\begin{theorem}[Abel--Jacobi]
雅可比簇 $J(X)$ 是阿贝尔簇 $\Pic^{0}(X)$ 的底层复环面，解析态射
\[
  \varphi:X\to J(X)\cong\Pic^{0}(X)
\]
正是把 $x\in X$ 送到 $\cO_X([x]-[x_0])$ 在 $\Pic^{0}(X)$ 中的同构类的代数态射。
\end{theorem}

由 $\varphi$ 诱导的同态
\[
  \pi_1(X,x_0)\to\pi_1(J(X),0)
\]
是商映射
\[
  \pi_1(X,x_0)\twoheadrightarrow
  \pi_1(X,x_0)/[\pi_1(X,x_0),\pi_1(X,x_0)]=\Lambda.
\]
所以对每个特征标 $\chi:\pi_1(X,x_0)\to\C^{\times}$，存在唯一的特征标
$\Aut_{\chi}:\pi_1(\Pic^{0}(X),0)\to\C^{\times}$，使 $\chi=\Aut_{\chi}\circ\varphi_{*}$。

流形上的局部系统（$\C$-向量空间）无非是该流形基本群的有限维 $\C$-表示。因此我们用解析方法证明了：

\begin{theorem}
对 $X$ 上任意秩 $1$ 的局部系统 $E$，存在 $\Pic^{0}(X)$ 上唯一的秩 $1$ 局部系统 $\Aut_E$（在原点刚性化），使 $\varphi^{*}\Aut_E=E$。此外，若 $m:\Pic^{0}(X)\times\Pic^{0}(X)\to\Pic^{0}(X)$ 是群运算，则有典范同构
\[
  m^{*}\Aut_E\cong\Aut_E\boxtimes\Aut_E.
\]
\end{theorem}

\srcpage{6}
\begin{proof}[定理 1.2 的代数证明]
对每个整数 $d\ge 0$，对称群 $\mathfrak{S}_d$ 以显然的方式作用在积
\[
  X^{d}=\overbrace{X\times\cdots\times X}^{d}
\]
上。

商概形
\[
  X^{(d)}=X^{d}/\mathfrak{S}_d
\]
连通、射影、光滑，维数为 $d$；它是 $X$ 上有效除子 $D=\sum_x d_x[x]$（$d_x\ge 0$）的模空间，这些除子的次数为 $\sum_x d_x=d$。态射 $\varphi$ 诱导态射
\[
  \varphi^{(d)}:X^{(d)}\to\Pic^{0}(X),\quad D\mapsto\cO_X(D-d[x_0]).
\]
只要 $d\ge 2g-1$，$\varphi^{(d)}$ 就是秩 $d-g$ 的射影丛，因此其纤维单连通。

对 $X$ 上每个秩 $n$ 的局部系统 $E$ 与每个整数 $d\ge 0$，外张量积 $E^{\boxtimes d}$ 是 $X^{d}$ 上秩 $n^{d}$ 的局部系统，带有 $\mathfrak{S}_d$ 的作用，该作用提升 $X^{d}$ 上的作用。记商态射为 $r:X^{d}\twoheadrightarrow X^{d}/\mathfrak{S}_d=X^{(d)}$，则 $\C$-向量空间的可构造层 $r_{*}E^{\boxtimes d}$ 带有 $\mathfrak{S}_d$ 的作用。记
\[
  E^{(d)}=(r_{*}E^{\boxtimes d})^{\mathfrak{S}_d}\subset E^{\boxtimes d}
\]
为该作用的不动向量子层。可以验证，$E^{(d)}$ 在 $D\in X^{(d)}$ 处的纤维等于
\[
  \tag*{(1.3)} (E^{(d)})_D=\bigotimes_x\Sym^{d_x}E_x.
\]

若现在 $E$ 是 $X$ 上秩 $1$ 的局部系统，则 $E^{(d)}$ 实际上是 $X^{(d)}$ 上秩 $1$ 的局部系统。对每个 $d\ge 2g-1$，$E^{(d)}$ 因此在态射 $\varphi^{(d)}:X^{(d)}\to\Pic^{0}(X)$ 的单连通纤维上常值，并下降为 $\Pic^{0}(X)$ 上的局部系统，正是 $\Aut_E$。
\end{proof}

这是处处非分歧的 Drinfeld--Langlands 对应中的 $\mathrm{GL}(1)$ 情形。

% ============================================================
%  SY02-46 · §2 主定理 · 印刷页 6–8（底本 p6 下段起，p7、p8）
%  条目：Definition 2.1、Theorem 2.2（FGV 主定理）
% ============================================================
\section{主定理}\srcpage{6}

设 $n\ge 2$ 是整数。$\Pic(X)$ 的自然推广，更确切地说是概形 $\Pic(X)$ 在 $\Gm$ 的平凡作用下的商叠的自然推广，是 Riemann 面 $X$ 上秩 $n$ 的向量丛的叠 $\Fib_n(X)$。这个叠是代数的，光滑且纯维数为 $n^2(g-1)$，并且是有限型开集的递增并。若 $g\ge 2$，它含一个稠密开集，该开集是 $X$ 上稳定丛的拟射影模概形上的 $\Gm$-gerbe。$\Fib_n(X)$ 的连通分支 $\Fib_n^d(X)$ 由泛向量丛的次数 $d$ 划分。

\srcpage{7}

$X$ 上秩 $n$ 的向量丛 $\cL$ 也可以看作秩 $n$ 的局部自由 $\cO_X$-模。$\cL$ 在点 $x\in X$ 处的初等下修改，指 $X$ 上秩 $n$ 的向量丛 $\cL'$，它作为子 $\cO_X$-模含于 $\cL$ 中，使得凝聚 $\cO_X$-模 $\cL/\cL'$ 是集中于 $x$、长度为 $1$ 的摩天层。因此 $\cL'$ 的次数等于
\[
  \deg(\cL') = \deg(\cL) - 1.
\]

Hecke 叠是三元组 $(x,\cL,\cL'\subset\cL)$ 的模叠 $\Hck_n(X)$，其中 $x\in X$，$\cL\in\Fib_n(X)$，而 $\cL'\subset\cL$ 是 $\cL$ 在 $x$ 处的初等下修改。

Hecke 对应是图表
\[
  \begin{tikzcd}[column sep=1.8em, row sep=1.8em]
    & \Hck_n(X) \arrow[dl, "{(q,p')}"'] \arrow[dr, "p"] & \\
    X\times\Fib_n(X) & & \Fib_n(X)
  \end{tikzcd}
\]

其中投影 $(q,p')$ 与 $p$ 由
\[
  (q,p')(x,\cL,\cL'\subset\cL) = (x,\cL')
\]
与
\[
  p(x,\cL,\cL'\subset\cL) = \cL
\]
给出，并且二者分别是相对维数 $n-1$ 与 $n$ 的可表、射影、光滑态射。

这个对应在导出范畴 $\Dbc(\Fib_n(X))$ 上诱导一个算子，称为 Hecke 算子。它把 $K\in\Dbc(\Fib_n(X))$ 送到
\[
  H(K) = (q,p')_* p^*K[n-1] \in \Dbc(X\times\Fib_n(X)).
\]

可以「迭代」$H$，特别地得到算子
\[
  H\circ H : \Dbc(\Fib_n(X)) \to \Dbc(X\times X\times\Fib_n(X)),
\]
它在对角线的余集 $(X\times X-\Delta_X)\times\Fib_n(X)$ 上的限制，关于置换 $X$ 的两个副本的作用是 $S_2$-等变的。

\begin{definition}
设 $E$ 是 $X$ 上秩 $n$ 的不可约局部系统。$\Fib_n(X)$ 上的不可约反常层 $K$ 称为相对 $E$ 具有 Hecke 性质，若存在同构
\[
  H(K)\cong E\boxtimes K,
\]
使得诱导同构
\[
  (H\circ H)(K)\cong E\boxtimes E\boxtimes K
\]
在开集 $(X\times X-\Delta_X)\times\Fib_n(X)$ 上的限制是 $S_2$-等变的。
\end{definition}

\srcpage{8}

\begin{theorem}[Drinfeld 于 $n=2$，{[Dr 2]}；一般情形由 Frenkel、Gaitsgory 与 Vilonen 给出，{[F-G-V 2]} 与 {[Ga]}]
对 $X$ 上任一秩 $n$ 的不可约局部系统 $E$，存在 $\Fib_n(X)$ 上的反常层 $\Aut_E$。它到每个连通分支 $\Fib_n^d(X)$ 的限制不可约，并且相对 $E$ 具有 Hecke 性质。
\end{theorem}

定理的证明分两步。

首先，在偶 $(L,s)$ 的模空间 $\Fib'_n(X)$ 上构造可构造层的复形 $\Aut'{}_E$。这里的偶 $(L,s)$ 由 $X$ 上秩 $n$ 的向量丛 $L$ 与一个模数乘的非零截面
\[
  s : (\OmX)^{\otimes(n-1)} \hookrightarrow L
\]
组成。

$\mathrm{GL}(n)$ 尖点自守形式的 Fourier 级数展开，以及关于该级数系数的 Shintani 公式，启发了这个构造；其中还考虑到 Grothendieck 的函数–层字典。但它不是简单的照搬。

然后证明：在 $\Fib'_n(X)$ 的一个大的开集上，$\Aut'{}_E$ 是反常层；它沿遗忘截面 $s$ 的态射 $\Fib'_n(X)\to\Fib_n(X)$ 下降，成为 $\Fib_n(X)$ 上具有所需性质的反常层 $\Aut_E$。当然，全部困难都在第二步。

% ============================================================
%  SY02-46 · §3 反常层 ℒ^d_E · 印刷页 8–9（底本 p8 下段起，p9）
%  条目：Lemma 3.1、Definition 3.2
% ============================================================
\section{反常层 $\cL^d_E$}\srcpage{8}

对每个整数 $d\ge 0$，设 $\Coh^d_0(X)$ 是一般秩 $0$ 的凝聚 $\cO_X$-模 $\cM$ 的叠，也就是说 $\cM$ 具有有限支撑，且长度为 $\dim H^0(X,\cM)=d$。记 $\Coh_0(X)$ 为诸 $\Coh^d_0(X)$ 的不交并。

叠 $\Coh^d_0(X)$ 的定义对任意拟射影代数曲线 $X$ 有意义；特别地，可以取叠 $\Coh^d_0(\A^1)$。一般秩 $0$、长度为 $d$ 的凝聚 $\cO_{\A^1}$-模，等价于一个 $d$ 维向量空间连同其一个自同态。因此 $\Coh^d_0(\A^1)$ 就是 $d\times d$ 方阵的仿射空间在 $\mathrm{GL}(d)$ 共轭作用下的商叠 $[\mathfrak{gl}(d)/\mathrm{GL}(d)]$。由于任意光滑代数曲线 $X$ 在 étale 拓扑下局部同构于仿射直线 $\A^1$，$\Coh^d_0(X)$ 在 étale 拓扑下局部同构于 $[\mathfrak{gl}(d)/\mathrm{GL}(d)]$。

$X$ 上每个有效除子 $D$ 定义挠 $\cO_X$-模
\[
  \cO_{X,D} = \cO_X/\cO_X(-D),
\]
其长度等于 $D$ 的次数。并且对每个整数 $d\ge 0$，有态射
\[
  X^{(d)} \to \Coh^d_0(X),\quad D\mapsto \cO_{X,D},
\]
\srcpage{9}
它光滑，相对维数为 $d$。

任意 $\cM\in\Coh_0(X)$ 都同构于 $\bigoplus_{i\in I}\cO_{X,D_i}$，其中 $(D_i)_{i\in I}$ 是 $X$ 上有效除子的有限族。并且对每个整数 $d\ge 0$，有行列式态射
\[
  \det : \Coh^d_0(X) \to X^{(d)},
\]
它把 $\cM\cong\bigoplus_{i\in I}\cO_{X,D_i}$ 送到 $\sum_{i\in I}D_i$。记 $X^{(d)}_{\rss}$ 为 $X^{(d)}$ 中由无重除子 $D$ 构成的开集，即对任意 $x$ 有 $d_x=0$ 或 $1$；记 $\Coh^d_{0,\rss}(X)=\det^{-1}(X^{(d)}_{\rss})$ 为 $\Coh^d_0(X)$ 中对应的开集。

设 $E$ 是 $X$ 上秩 $n$ 的局部系统。关于 $E^{(d)}$ 在 $D\in X^{(d)}$ 处纤维的公式 (1.3) 表明，该纤维的秩随 $D$ 变化；但在开集 $X^{(d)}_{\rss}$ 之上，这个秩是常值，且 $E^{(d)}$ 的限制是秩 $n^d$ 的局部系统。

\begin{lemma}
对每个整数 $d\ge 0$，复形 $E^{(d)}[d]\in\Dbc(X^{(d)})$ 由可构造层 $E^{(d)}$ 放在次数 $-d$ 处构成。它是反常层，并且是其到 $X^{(d)}_{\rss}$ 上限制的中间延拓。反常层 $E^{(d)}[d]$ 在 $E$ 不可约（相应地，半单）时不可约（相应地，半单）。
\end{lemma}

把可构造层 $\det^* E^{(d)}$ 限制到 $\Coh^d_0(X)$ 的稠密开集 $\Coh^d_{0,\rss}(X)$ 上，得
\[
  \cL^d_{E,\rss} = \det^* E^{(d)}|_{\Coh^d_{0,\rss}(X)},
\]
它也是秩 $n^d$ 的局部系统，并且在 $E$ 不可约（相应地，半单）时不可约（相应地，半单）。

\begin{definition}
对每个整数 $d\ge 0$，与 $X$ 上局部系统 $E$ 相伴的反常层 $\cL^d_E$，就定义为 $\cL^d_{E,\rss}$ 到整个 $\Coh^d_0(X)$ 的中间延拓。
\end{definition}

当然，若 $E$ 不可约（相应地，半单），则 $\cL^d_E$ 不可约（相应地，半单）。此外可以验证，$\cL^d_E$ 沿态射 $X^{(d)}\to\Coh^d_0$，$D\mapsto\cO_{X,D}$ 的 $d$-平移拉回，正是反常层 $E^{(d)}[d]$。

\section{复形 ${\Aut}'_E$}\srcpage{9}

此后假设 $X$ 的亏格 $g\ge 2$。当亏格为 $0$ 或 $1$ 时，$\mathrm{GL}(n)$（$n\ge 2$）的处处非分歧
Drinfeld--Langlands 对应是空的，因为 $X$ 上没有秩 $\ge 2$ 的不可约局部系统。

本节打算构造一个候选 ${\Aut}'_E$，作为 $\Aut_E$ 到代数叠 $\Fib'_n(X)$ 上的限制。该叠的对象是
$X$ 上秩 $n$ 的向量丛 $\cL$，附带一个截面 $s:(\OmX)^{\otimes n-1}\to\cL$（模数乘）。

考虑由三元组 $(\cL,\cL_\bullet,s_\bullet)$ 构成的代数叠 $\cU_n$，其中 $\cL_n$ 是秩 $n$ 的向量丛，
\[
  \cL_\bullet=((0)=\cL_0\subset\cL_1\subset\cdots\subset\cL_n=\cL)
\]
\srcpage{10}
是子向量丛的完全旗，且 $s_\bullet$ 是一族线丛同构
\[
  s_j:(\OmX)^{\otimes j-1}\xrightarrow{\ \sim\ }\cL_{n-j+1}/\cL_{n-j},\qquad j=1,\dots,n,
\]
这族同构在模数乘下等同：对每个 $t\in\Gm$，有
$(s_j)_{j=1,\dots,n}\sim(ts_j)_{j=1,\dots,n}$。

有态射
\[
  \begin{tikzcd}
    \cU_n \arrow[r,"g_n"] \arrow[d,swap,"f_n"] & \A^1 \\
    \Fib'^{n(n-1)(g-1)}_n(X) &
  \end{tikzcd}
\]

其中 $f_n$ 把三元组 $(\cL,\cL_\bullet,s_\bullet)$ 送到向量丛 $\cL$，并带上由下式诱导的截面（模数乘）：
\[
  s_n:(\OmX)^{\otimes n-1}\xrightarrow{\ \sim\ }\cL_1\subset\cL_n=\cL.
\]

而 $g_n$ 把同一个三元组送到标量，即下列扩张类之和：
\[
  0\to\cL_j/\cL_{j-1}\to\cL_{j+1}/\cL_{j-1}\to\cL_{j+1}/\cL_j\to 0
\]
这些扩张类落在
\[
  \begin{aligned}
    \Ext^1_{\cO_X}(\cL_{j+1}/\cL_j,\cL_j/\cL_{j-1})
    &\cong\Ext^1_{\cO_X}\bigl((\OmX)^{\otimes n-j-1},(\OmX)^{\otimes n-j}\bigr)\\
    &\cong\Ext^1_{\cO_X}(\cO_X,\OmX)\cong\C.
  \end{aligned}
\]

若让乘法群 $\Gm$ 作用在 $\cU_n$ 上，方式为
\[
  t\cdot(\cL,\cL_\bullet,s_\bullet)=(\cL,\cL_\bullet,t\cdot s_\bullet),
\]
其中 $t\cdot s_\bullet=(t^{n-j}s_j)_{j=1,\dots,n}$；并让 $\Gm$ 以数乘作用在 $\A^1$ 上，则态射
$g_n$ 是 $\Gm$-等变的。对该作用取商，得到态射
\[
  \begin{tikzcd}
    \overline{\cU}_n \arrow[r,"\overline g_n"] \arrow[d,swap,"\overline f_n"]
      & {[\A^1/\Gm]=\cA} \\
    \Fib'^{n(n-1)(g-1)}_n(X) &
  \end{tikzcd}
\]

于是可以构造可构造层的复形
\[
  L_n=(\overline f_n)_!(\overline g_n)^*\Psi[\dim\overline{\cU}_n],
\]
其中 $\Psi\in\Dbc(\cA)$ 是第 0 节定义的复形。

\srcpage{11}
对每个整数 $d$，有一个对应
\[
  \begin{tikzcd}
    & \Mod'^d_n(X) \arrow[r,"q"] \arrow[dl,swap,"p'"] \arrow[dr,"p"] & \Coh^d_0(X) \\
    \Fib'^{n(n-1)(g-1)}_n(X) & & \Fib'^{n(n-1)(g-1)+d}_n(X)
  \end{tikzcd}
\]
其中 $\Mod'^d_n(X)$ 是图表
\[
  (\OmX)^{\otimes n-1}\hookrightarrow\cL'\subset\cL
\]
的叠。这样的图表由两个向量丛与两个 $\cO_X$-线性单射组成：两个向量丛的秩均为 $n$，
次数分别为 $n(n-1)(g-1)$ 与 $n(n-1)(g-1)+d$；两个单射为
$(\OmX)^{\otimes n-1}\hookrightarrow\cL'$（模数乘）与 $\cL'\hookrightarrow\cL$。这里
$p$、$p'$、$q$ 分别把这样的图表送到 $(\OmX)^{\otimes n-1}\hookrightarrow\cL'$（模数乘）、
$(\OmX)^{\otimes n-1}\hookrightarrow\cL$（模数乘）与 $\cL/\cL'$。

若 $E$ 是 $X$ 上的局部系统，则定义函子\yizhu{下文函子定义右端写作 $\Fib'_i$，下标 $i$ 与本节其余
各处的 $n$ 不一致，按上下文应为 $n$。}{原文笔误，I（p11 上下文推断，无外部依据）}
\[
  M'^d_{n,E}:\Dbc\bigl(\Fib'^{n(n-1)(g-1)}_n(X)\bigr)\to\Dbc\bigl(\Fib'^{n(n-1)(g-1)+d}_i(X)\bigr)
\]
为
\[
  M'^d_{n,E}(K)=p_*p'^*(K\otimes q^*\cL^d_E)[nd],
\]
其中 $\cL^d_E$ 是第 3 节构造的 $\Coh^d_0(X)$ 上的反常层。

\begin{definition}
  若 $E$ 是 $X$ 上的局部系统，则在 $\Fib'_n(X)$ 上定义复形 ${\Aut}'_E$ 为：当 $d\ge 0$ 时
  \[
    {\Aut}'_E|_{\Fib'^{n(n-1)(g-1)+d}_n(X)}=M'^d_{n,E}(L_n),
  \]
  当 $d<0$ 时
  \[
    {\Aut}'_E|_{\Fib'^{n(n-1)(g-1)+d}_n(X)}=(0).
  \]
\end{definition}

\section{用 Fourier 变换构造}\srcpage{11}

对每个 $i=1,\dots,n$，设 $\Coh_i(X)$ 为一般秩为 $i$ 的凝聚 $\cO_X$-模 $\cM_i$ 构成的代数叠。
设 ${\Coh'}_i(X)$ 为一般秩为 $i$ 的凝聚 $\cO_X$-模 $\cM_i$ 再附上一个模数乘的
$\cO_X$-线性单射 $s_i\colon(\OmX)^{\otimes i-1}\hookrightarrow\cM_i$ 所得的代数叠。
这两个叠以 $\Fib_i(X)$ 与 ${\Fib'}_i(X)$ 为稠密开集。
与后两个叠一样，它们的连通分支 $\Coh^d_i(X)$ 与 ${\Coh'}^d_i(X)$ 按 $\cM_i$ 的次数 $d$ 划分。

\srcpage{12}
\begin{lemma}
存在常数 $c(g,n)$，使得以下结论成立。
对每个整数 $d\ge c(g,n)$ 与每个 $L\in\Fib^d_n(X)$，只要
$\Hom_{\cO_X}\big(L,(\OmX)^{\otimes n+1}\big)\ne(0)$，$L$ 就是非常不稳定的。
也就是说，存在非平凡的直和分解 $L=L_1\oplus L_2$，使 $\Ext^1_{\cO_X}(L_1,L_2)=(0)$。%
\hfill$\square$
\end{lemma}

对每个 $i=0,\dots,n$，设 $\cC_i\subset\Coh_i(X)$ 为由满足以下两条的 $\cM_i$ 构成的开集：
其次数 $\ge c(n,g)+i(i-1)(g-1)$，且 $\Hom_{\cO_X}\big(\cM_i,(\OmX)^{\otimes n+1}\big)=(0)$。
于是特别地 $\Hom_{\cO_X}\big(\cM_i,(\OmX)^{\otimes i}\big)=(0)$；
并且由 Serre 对偶，$\Ext^1_{\cO_X}\big((\OmX)^{\otimes i-1},\cM_i\big)=(0)$。
设 $\cV_i$ 为二元组 $(\cM_i,s_i)$ 的代数叠，其中 $\cM_i\in\cC_i$，而
$s_i\in\Hom_{\cO_X}\big((\OmX)^{\otimes i-1},\cM_i\big)$。
设 $\cV^\vee_i$ 为满足 $\cM_i\in\cC_i$ 的凝聚 $\cO_X$-模扩张
\[
  0\to(\OmX)^{\otimes i}\to\cM_{i+1}\to\cM_i\to 0
\]
的代数叠。有自然的投影 $\pi_i\colon\cV_i\to\cC_i$ 与 $\pi^\vee_i\colon\cV^\vee_i\to\cC_i$，
它们是对偶的向量丛。

引入开集
\[
  \cV{}^\circ_i=\big\{(\cM_i,s_i)\ \big|\ s_i\ \text{是单射}\big\}\subset\cV_i
\]
与
\[
  \cV{}^{\vee\circ}_i=\big\{(0\to(\OmX)^{\otimes i}\to\cM_{i+1}\to\cM_i\to 0)\ \big|\
  \cM_{i+1}\in\cC_{i+1}\big\}\subset\cV^\vee_i .
\]
当然，有典范同构
\[
  \cV{}^{\vee\circ}_i\xrightarrow{\ \sim\ }\cV{}^\circ_{i+1},
\]
它把 $0\to(\OmX)^{\otimes i}\to\cM_{i+1}\to\cM_i\to 0$ 送到单射
$(\OmX)^{\otimes i}\to\cM_{i+1}$。

乘法群按数乘作用在向量丛 $\cV_i$ 与 $\cV^\vee_i$ 上；取商即得基本图
\begin{center}
\footnotesize
\begin{tikzcd}[column sep=1.55em, row sep=1.7em]
  & \bar\cV{}^\circ_n
      \arrow[dl]
  & \bar\cV{}^{\vee\circ}_{n-1}
      \arrow[l, "\sim"']
      \arrow[r, hook]
  & \bar\cV{}^\vee_{n-1}
      \arrow[dr]
  & & \bar\cV_{n-1}
      \arrow[dl]
  & \bar\cV{}^\circ_{n-1}
      \arrow[l, hook']
  & \bar\cV{}^{\vee\circ}_{n-2}
      \arrow[l, "\sim"']
      \arrow[r, hook]
  & \cdots
  & & & \\
  \cC_n & & & & \cC_{n-1} & & & & & & & \\
  \cdots
  & \bar\cV{}^\vee_2
      \arrow[dr]
  & & \bar\cV_2
      \arrow[dl]
  & \bar\cV{}^\circ_2
      \arrow[l, hook']
  & \bar\cV{}^{\vee\circ}_1
      \arrow[l, "\sim"']
      \arrow[r, hook]
  & \bar\cV{}^\vee_1
      \arrow[dr]
  & & \bar\cV_1
      \arrow[dl]
  & \bar\cV{}^\circ_1
      \arrow[l, hook']
  & \bar\cV{}^{\vee\circ}_0
      \arrow[l, "\sim"']
      \arrow[dr]
  & \\
  & & \cC_2 & & & & & \cC_1 & & & & \cC_0
\end{tikzcd}
\end{center}
其中斜箭头是投影
$\bar\pi{}^\circ_n\colon\bar\cV{}^\circ_n=[\cV^\circ_n/\Gm]\to\cC_n$，
$\bar\pi_i\colon\bar\cV_i=[\cV_i/\Gm]\to\cC_i$，
$\bar\pi{}^\vee_i\colon\bar\cV{}^\vee_i=[\cV^\vee_i/\Gm]\to\cC_i$，
$\bar\pi{}^{\vee\circ}_0\colon\bar\cV{}^{\vee\circ}_0=[\cV^{\vee\circ}_0/\Gm]\to\cC_0$。
水平箭头一方面是开浸入
$j_i\colon\bar\cV{}^\circ_i\hookrightarrow\bar\cV_i$ 与
$j^\vee_i\colon\bar\cV{}^{\vee\circ}_i\hookrightarrow\bar\cV{}^\vee_i$，
它们由包含 $\cV^\circ_i\subset\cV_i$ 与 $\cV^{\vee\circ}_i\subset\cV^\vee_i$ 诱导；
另一方面是同构
$\iota_i\colon\bar\cV{}^{\vee\circ}_i\xrightarrow{\ \sim\ }\bar\cV{}^\circ_{i+1}$，
它们由上面的同构 $\cV{}^{\vee\circ}_i\xrightarrow{\ \sim\ }\cV{}^\circ_{i+1}$ 诱导。

\srcpage{13}
于是对 $X$ 上每个秩为 $n$ 的局部系统 $E$，定义复形
$K_{E,i}\in\Dbc(\bar\cV{}^\circ_i)$，$i=1,\dots,n$，如下：
\[
  K_{E,1}=(\iota_0)_*\,(\bar\pi{}^{\vee\circ}_0)^*\,\cL_E,
\]
其中 $\cL_E$ 是 $\cC_0$ 上的复形：它在 $\cC_0$ 的每个连通分支 $\cC^d_0$ 上的限制，
是 3.2 中定义的反常层 $\cL^d_E$ 平移 $d-1$；并由递推关系
\[
  K_{E,i+1}=(\iota_i)_*\,(j^\vee_i)^*\,\Four_{\bar\cV_i/\cC_i}\big((j_i)_!K_{E,i}\big),
  \qquad i=1,\dots,n-1
\]
定义，其中 $\Four_{\bar\cV_i/\cC_i}\colon\Dbc(\bar\cV_i)\to\Dbc(\bar\cV{}^\vee_i)$ 是向量丛
$\cV_i\to\cC_i$ 的齐次 Fourier 变换，其定义见第~0~节。

叠 $\bar\cV{}^\circ_i$ 与 ${\Fib'}_i(X)$ 都是 ${\Coh'}_i(X)$ 的开集，因此可以考虑它们的交。

\begin{lemma}
设 $E$ 是 $X$ 上秩为 $n$ 的局部系统。
对每个 $i=1,\dots,n$ 与每个整数 $d\ge 0$，$K_{E,i}$ 在
${\Coh'}^{i(i-1)(g-1)+d}_i(X)$ 的开集
$\bar\cV{}^\circ_i\cap{\Fib'}^{i(i-1)(g-1)+d}_i(X)$ 上的限制等于
\[
  M'^{d}_{i,E}(L_i),
\]
这里采用第 4 节的记号。

特别地，对每个整数 $d\ge 0$，在 ${\Coh'}^{n(n-1)(g-1)+d}_i(X)$ 的开集
$\bar\cV{}^\circ_n\cap{\Fib'}^{n(n-1)(g-1)+d}_i(X)$ 上，$K_{E,n}$ 与 $\Aut'{}_E$ 的限制相合。%
\yizhu{末句两处下标写的都是 $i$，与句首的 $K_{E,n}$、$\Aut'{}_E$ 以及本引理首段不合，
似应为 $n$；此处照原文翻译。}{原文笔误，I（据 p13 上下文，无外部依据）}\hfill$\square$
\end{lemma}

\section{消没定理}\srcpage{13}

设 $i\ge 1$ 与 $d\ge 0$ 为整数。考虑图表
\[
\begin{tikzcd}
 & \Mod^d_i(X) \arrow[dl,"p'"] \arrow[r,"q"] \arrow[dr,"p"] & \Coh^d_0(X) \\
\Fib_i(X) & & \Fib_i(X)
\end{tikzcd}
\]

其中 $\Mod^d_i(X)$ 是由对 $(\cL,\cL'\subset\cL)$ 构成的叠。这里 $\cL$ 是秩 $i$ 的向量丛，$\cL'\subset\cL$ 是余长 $d$ 的下修改。投影 $p'$ 与 $p$ 分别把 $(\cL,\cL'\subset\cL)$ 送到 $\cL'$ 与 $\cL$。$q$ 则把它送到凝聚 $\cO_X$-模 $\cL/\cL'$，后者一般秩为 $0$、长度为 $d$。

对 $X$ 上任一局部系统 $E$，第 4 节的函子 $M'^d_{i,E}$ 有如下变体
\[
M^d_{i,E}:\Dbc(\Fib_i(X))\to\Dbc(\Fib_i(X)),\qquad
M^d_{i,E}(K)=p_*p'^*(K\otimes q^*\cL^d_E)[id],
\]
\srcpage{14}
其中 $\cL^d_E$ 是 $\Coh^d_0(X)$ 上由 3.2 定义的反常层。

\begin{theorem}[Gaitsgory, {[Ga]}]
设 $E$ 是秩 $n$ 的不可约局部系统。则对一切 $i=1,\dots,n-1$ 与一切整数 $d>in(2g-2)$\yizhu{原文作 $in(2g-2)$，其中 $in$ 是 $i$ 与 $n$ 之积，即 $d>i\,n(2g-2)$；注 6.2 中 $i=1$ 的情形正写作 $n(2g-2)$，可互证。}{术语背景，A（原文注 6.2 互证）}，函子 $M^d_{i,E}$ 恒为零。
\end{theorem}

\begin{remark}
这个定理的证明很长，技术性很强。它远远超出了本报告的范围。然而我们指出，当 $i=1$ 时，这一陈述不过是 Deligne 一个结果的重新表述（见 [Dr 2]），该结果说
\[
(\varphi^{(d)})_*E^{(d)}=0
\]
对 $X$ 上任一秩 $n>1$ 的不可约局部系统与任一整数 $d>n(2g-2)$ 成立，其中 $\varphi^{(d)}:X^{(d)}\to\Pic^0(X)$ 是第 1 节中考虑的射影丛。
\end{remark}

Frenkel、Gaitsgory 与 Vilonen 用定理 6.1 证明下述命题。

\begin{proposition}
设 $E$ 是秩 $n$ 的不可约局部系统。则对每个 $i=1,\dots,n-1$，典范态射
\[
(j_i)_!K_{E,i}\to(j_i)_*K_{E,i}
\]
是同构。
\end{proposition}

证明思路：在 $\cC_i$ 的开集 $\Fib_i(X)\cap\cC_i$ 上，$\cV^\circ_i\subset\cV_i$ 是向量丛 $\cV_i\to\cC_i$ 中零截面的补集。由引理 (0.3)，要证明该开集上的断言，只需证明 $(\overline{\pi}^\circ_i)_!K_{E,i}=(0)$，其中 $\overline{\pi}^\circ_i:\overline{\cV}^\circ_i\to\cC_i$ 是投影；或等价地，证明关于投影 $\Pv(\cV_i)\to\cC_i$ 的同一等式，因为这两个投影在 $\Fib_i(X)\cap\cC_i$ 上重合。而由引理 5.2，对每个整数 $d\ge 0$，有
\[
(\overline{\pi}^\circ_i)_!K_{E,i}|(\cC_i\cap\Fib_i^{i(i-1)(g-1)+d}(X))
=M^d_{i,E}(\pi'_*L_i)|(\cC_i\cap\Fib_i^{i(i-1)(g-1)+d}(X))
\]
其中 $\pi':\Fib'_i(X)\to\Fib_i(X)$ 是遗忘截面 $s:(\OmX)^{\otimes n-1}\to\cL$ 的态射，于是由定理 6.1 得结论。

由命题 6.3，若 $K_{E,i}$ 是不可约反常层，则 $(j_i)_!K_{E,i}$ 也是，因为此时它是 $K_{E,i}$ 的中间延拓。由于齐次 Fourier 变换保持反常性与不可约性，对 $i$ 作归纳可知：$K_{E,i}$ 在 $\overline{\cV}^\circ_i$ 的每个连通分支上的限制是不可约反常层。特别地，我们已经证明：

\srcpage{15}
\begin{corollary}
设 $E$ 是秩 $n$ 的不可约局部系统。则 $K_{E,n}$ 在 $\overline{\cV}^\circ_n$ 的每个连通分支上的限制是不可约反常层。
\end{corollary}

\section{下降}\srcpage{15}

我们先回忆两个一般结果。对任意概形（或代数叠）$S$ 与任意 $K \in \Dbc(S)$，记
\[
  \chi_K \colon S \to \Z
\]
为 $K$ 的 Euler--Poincar\'e 示性函数，它由
\[
  \chi_K(s) = \sum_i (-1)^i \dim_{\C} H^i(K_s), \qquad \forall s \in S,
\]
定义，其中 $K_s$ 是 $K$ 在 $s$ 处的纤维。

\begin{lemma}[Deligne, {[Il]} 推论 2.10]
设 $f \colon Y \to Z$ 是概形的紧合态射，甚至是代数叠的可表紧合态射。那么，若
$K, K' \in \Dbc(Y)$ 在 $Y$ 上局部同构，则复形 $f_*K$ 与 $f_*K'$ 在 $Z$ 上有相同的
Euler--Poincar\'e 示性函数 $\chi_{f_*K} = \chi_{f_*K'}$。
\end{lemma}

\begin{lemma}
设 $S$ 是概形，甚至是代数叠；设 $V \to S$ 是常秩 $r$ 的向量丛（$r$ 有限），它定义射影丛
$\pi \colon P = \Pv(V) \to S$；设 $K$ 是 $P$ 上的不可约反常层。那么，存在 $S$ 上的反常层
（必然不可约）$L$，使得 $K$ 同构于 $\pi^*(L)[r-1]$，当且仅当
Euler--Poincar\'e 示性函数 $\chi_K$ 沿 $\pi$ 的纤维为常数。
\end{lemma}

\begin{proof}
「仅当」部分平凡。因此设 $\chi_K$ 沿 $\pi$ 的纤维为常数。由不可约反常层的结构定理，有
$K \cong i_* j_{!*} E[\dim Q]$，其中 $i \colon Q \hookrightarrow P$ 是 $K$ 的不可约闭支撑，
$j \colon Q^\circ \hookrightarrow Q$ 是稠密开集，$E$ 是 $Q^\circ$ 上的不可约局部系统。

首先验证 $Q = \pi^{-1}(T)$，其中 $T$ 是 $S$ 的不可约闭子集。其次验证可取 $Q^\circ$ 形如
$\pi^{-1}(T^\circ)$，其中 $T^\circ$ 是 $T$ 的稠密开集。最后验证 $E$ 来自 $T^\circ$ 上的
不可约局部系统 $F$，因为 $\pi$ 的纤维单连通。

所求的 $S$ 上反常层 $L$，就是 $F[\dim T]$ 到 $T$ 的中间延拓在整个 $S$ 上的零延拓。
\end{proof}

现在研究 $\Aut'_E$ 到 $\Fib_n(X)$ 上反常层 $\Aut_E$ 的下降。为此，考虑 Drinfeld 引入的叠
$\cY^d_n$，它分类满足下述条件的二元组 $(\cL, s^\bullet)$：
\srcpage{16}$\cL$ 是次数为 $n(n-1)(g-1)+d$ 的向量丛，而
$s^\bullet = (s^1, \dots, s^n)$ 是一列 $\cO_X$-模单射
\[
\begin{array}{rcl}
(\OmX)^{\otimes n-1} & \xhookrightarrow{s^1} & \cL \\
 & \cdots & \\
(\OmX)^{\otimes (n-1)+\cdots+(n-i)} & \xhookrightarrow{s^i} & \bigwedge^i \cL \\
 & \cdots & \\
(\OmX)^{\otimes \frac{n(n-1)}{2}} & \xhookrightarrow{s^n} & \bigwedge^n \cL
\end{array}
\]
它们满足 Pl\"ucker 关系，使得 $s^\bullet$ 在 $X$ 的一般点处给出 $\cL$ 纤维的一个
子空间完全旗。

叠 $\cY^d_n$ 是三元组 $(\cL, \cL_\bullet, s_\bullet)$ 的叠的一个
「部分紧化」，其中
\[
  \cL_\bullet = ((0) = \cL_0 \subset \cL_1 \subset \cdots \subset \cL_n)
\]
是 $\cL$ 的 $\cO_X$-子模旗，而 $s_\bullet$ 是一列 $\cO_X$-模同构
$s_j \colon (\OmX)^{\otimes j-1} \xrightarrow{\sim} \cL_{n-j+1}/\cL_{n-j}$，
$j = 1, \dots, n$。事实上，由
\[
\begin{aligned}
  s^i = s_n \otimes s_{n-1} \otimes \cdots \otimes s_{n-i+1} \colon {}
  & (\OmX)^{\otimes (n-1)+(n-2)+\cdots+(n-i)} \\
  & \xrightarrow{\sim}
  \cL_1 \otimes (\cL_2/\cL_1) \otimes \cdots \otimes
  (\cL_i/\cL_{i-1}) \cong \bigwedge^i \cL_i
\end{aligned}
\]
定义的态射
\[
  (\cL, \cL_\bullet, s_\bullet) \mapsto (\cL, s^\bullet)
\]
把这个三元组叠与开集 $\cY^{d\circ}_n \subset \cY^d_n$ 等同，该开集由满足
下述条件的 $(\cL, s^\bullet)$ 组成：当 $i = 1, \dots, n-1$ 时，单射
\[
  s^i \colon (\OmX)^{\otimes (n-1)+\cdots+(n-i)} \hookrightarrow \bigwedge^i \cL
\]
的余核无挠。

环面 $\GG_m^n$ 按下式作用在 $\cY^d_n$ 上：
\[
  (t_1, \dots, t_n) \cdot (\cL, (s^1, \dots, s^n)) =
  (\cL, (t_1 s^1, \dots, t_n s^n)).
\]

这个作用保持开集 $\cY^{d\circ}_n$，它在三元组叠上由
\[
  (t_1, \dots, t_n) \cdot (\cL, \cL_\bullet, (s_1, \dots, s_{n-1}, s_n)) =
  (\cL, \cL_\bullet,
  (t_n t_{n-1}^{-1} s_1, \dots, t_2 t_1^{-1} s_{n-1}, t_1 s_n))
\]
给出。

特别地，第 4 节引入的代数叠 $\cU_n$（相应地，
$\overline{\cU}_n = [\cU_n/\GG_m]$），可以等同于
$\cY^{0\circ}_n$ 对 $\GG_m$（相应地，$\GG_m^2$）作用的商。该作用经由嵌入
$\GG_m \hookrightarrow \GG_m^n$，$t \mapsto (t, t^2, \dots, t^n)$（相应地，
$\GG_m^2 \hookrightarrow \GG_m^n$，
$(t, t') \mapsto (t, t^2t', t^3t'^3, \dots, t^n t'^{\frac{n(n-1)}{2}})$）给出。

\srcpage{17}
\begin{proposition}[{[B-G]} 命题 1.2.2]
遗忘态射 $(\cL, s^\bullet) \to (\cL, s^1)$ 过商为代数叠的态射
\[
  \tilde{f} \colon \widetilde{\cY}^d_n := [\cY^d_n/\GG_m^n]
  \to {\Fib'}^{n(n-1)(g-1)+d}_n(X),
\]
它是可表且紧合的。
\end{proposition}

考虑第 4 节引入的 $\overline{\cU}_n = [\cY^{0\circ}_n/\GG_m^2]$ 上的复形
$(\overline{g}_n)^*\Psi$，并记 $\Phi \in \Dbc(\widetilde{\cY}^0_n)$ 为它沿复合
态射的紧支撑前推；该复合态射由商映射
$[\cY^{0\circ}_n/\GG_m^2] \to [\cY^{0\circ}_n/\GG_m^n]$ 与包含
$[\cY^{0\circ}_n/\GG_m^n] \subset [\cY^0_n/\GG_m^n] =
\widetilde{\cY}^d_n$ 组成。\yizhu{此处原文印作
$\widetilde{\cY}^d_n$。按上式末端的等号，$[\cY^0_n/\GG_m^n]$ 应等于
$\widetilde{\cY}^0_n$（见 p17 处的定义）；且 $\Phi$ 本身是
$\widetilde{\cY}^0_n$ 上的复形。故 $d$ 应为 $0$。译文照原文保留。}{原文笔误，I（原文 p15 定义互证）}

对 $X$ 上任意秩 $n$ 局部系统 $E$，有函子 $M^d_{n,E}$ 的一个变体
\[
  N^d_{n,E} \colon \Dbc(\widetilde{\cY}^0_n) \to \Dbc(\widetilde{\cY}^d_n),
  \quad K \mapsto p_*p'^*(K \otimes q^*\cL^d_E)[nd],
\]
它借助图表
\[
\begin{tikzcd}
 & \widetilde{\cZ}^d_n \arrow[dl, "p'"'] \arrow[dr, "p"] \arrow[r, "q"]
 & \Coh^d_0(X) \\
 \widetilde{\cY}^0_n & & \widetilde{\cY}^d_n
\end{tikzcd}
\]
定义。其中
\[
  \widetilde{\cZ}^d_n = \widetilde{\cY}^0_n
  \times_{\Fib_n(X),p'} \Mod^d_n(X)
\]
是三元组 $(\cL', s'^\bullet, \cL' \subset \cL)$ 的代数叠对
$\GG_m^n$ 的商，这里 $(\cL', s'^\bullet) \in \cY^0_n$ 且
$(\cL' \subset \cL) \in \Mod^d_i(X)$\yizhu{原文此处印作 $\Mod^d_i(X)$；按本页纤维积的定义与 $\cL'$ 的秩 $n$，下标应为 $n$。译文照原文保留。}{原文笔误，I（p17 上下文推断，无外部依据）}；而 $p$、$p'$、$q$ 由下述箭头诱导：
这些箭头把这样的三元组分别送到 $(\cL, s^\bullet)$、$(\cL', s'^\bullet)$ 与
$\cL/\cL'$，其中 $s^i$ 是 $s'^i$ 与包含
$\bigwedge^i \cL' \subset \bigwedge^i \cL$ 的复合。

\begin{lemma}
设 $E$ 是 $X$ 上秩 $n$ 的局部系统。对任意整数 $d \ge 0$，有
\[
  \Aut'_E\,|\,{\Fib'}^{n(n-1)(g-1)+d}_n = \tilde{f}_*N^d_{n,E}(\Phi).
\]
\end{lemma}

为证明 $\Aut'_E$ 沿遗忘截面的态射 $\Fib'_n(X) \to \Fib_n(X)$ 下降，Frenkel、Gaitsgory
与 Vilonen 用引理 7.4、命题 7.3 与引理 7.1 证明 $\Aut'_E$ 的
Euler--Poincar\'e 示性函数不依赖 $E$。然后他们在 $E$ 为秩 $n$ 平凡局部系统时计算这个
Euler--Poincar\'e 示性函数，最后用引理 7.2 得出结论。

\section*{参考文献}\srcpage{18}
\sloppy

\refitem{[B-B-D]}{A. Beilinson, J. Bernstein, P. Deligne -- Faisceaux pervers,
  \textit{Ast\'erisque} \textbf{100}, (1982).}

\refitem{[B-D]}{A. Beilinson, V. Drinfeld -- Quantization of Hitchin's Integrable System
  and Hecke Eigensheaves, pr\'epublication,
  http ://www.math.uchicago.edu/$\sim$benzvi/BD/hitchin.ps.gz.}

\refitem{[B-L]}{J. Bernstein, V. Lunts -- Equivariant Sheaves and Functors, Lecture Notes
  in Mathematics \textbf{1578}, Springer-verlag (1974).}

\refitem{[Br]}{J.-L. Brylinski -- Transformations canoniques, Dualit\'e projective, Th\'eorie
  de Lefschetz, Transformations de Fourier et Sommes trigonom\'etriques, dans
  \textit{G\'eom\'etrie et Analyse Microlocales}, \textit{Ast\'erisque} \textbf{140-141},
  (1986), 3-134.}

\refitem{[B-G]}{A. Braverman, D. Gaitsgory -- Geometric Eisenstein series, pr\'epublication,
  http ://arXiv.org/abs/math/9912097, (1999).}

\refitem{[Dr 1]}{V. G. Drinfeld -- Langlands' conjecture for GL(2) over functional fields,
  Proceedings of the International Congress of Mathematicians (Helsinki, 1978),
  \textit{Acad. Sci. Fennica}, Helsinki, (1980), 565-574.}

\refitem{[Dr 2]}{V. G. Drinfeld -- Two-dimensional $l$-adic representations of the fundamental
  group of a curve over a finite field and automorphic forms on GL(2),
  \textit{Amer. J. Math.} \textbf{105}, (1983), 85-114.}

\refitem{[Dr 3]}{V. G. Drinfeld -- Two-dimensional $l$-adic representations of the Galois group
  of a global field of characteristic $p$ and automorphic forms on GL(2),
  \textit{J. of Soviet Math.} \textbf{36}, (1987), 93-105.}

\refitem{[F-G-K-V]}{E. Frenkel, D. Gaitsgory, D. Kazhdan, K. Vilonen -- Geometric realization
  of Whittaker functions and the Langlands conjecture,
  \textit{J. Amer. Math. Soc.} \textbf{11}, (1998), 451-484.}

\refitem{[F-G-V 1]}{E. Frenkel, D. Gaitsgory, K. Vilonen -- Whittaker patterns in the geometry
  of moduli spaces of bundles on curves, \textit{Annals of Math.} \textbf{153}, (2001), 699-748.}

\refitem{[F-G-V 2]}{E. Frenkel, D. Gaitsgory, K. Vilonen -- On the geometric Langlands
  conjecture, \textit{J. Amer. Math. Soc.} \textbf{15}, (2002), 367-417.}

\refitem{[Ga]}{D. Gaitsgory -- On a vanishing conjecture appearing in the geometric Langlands
  correspondence, pr\'epublication, http ://arXiv.org/abs/math/0204081, (2002).}

\refitem{[Il]}{L. Illusie -- Th\'eorie de Brauer et caract\'eristiques d'Euler-Poincar\'e
  d'apr\`es Deligne, \textit{Ast\'erisque} \textbf{82-83}, (1981), 161-172.}

\srcpage{19}
\refitem{[La]}{L. Lafforgue -- Chtoucas de Drinfeld et correspondance de Langlands,
  \textit{Invent. Math.} \textbf{147}, (2002), 1-241.}

\refitem{[Lau 1]}{G. Laumon -- Correspondance de Langlands g\'eom\'etrique pour les corps de
  fonctions, \textit{Duke Math. J.} \textbf{54}, (1987), 309-359.}

\refitem{[Lau 2]}{G. Laumon -- Faisceaux automorphes pour GL$_n$ : la premi\`ere construction de
  Drinfeld, pr\'epublication, http ://arXiv.org/abs/math/9511004, (1995).}

\refitem{[Lau 3]}{G. Laumon -- Transformation de Fourier homog\`ene, en pr\'eparation.}

\refitem{[L-M]}{G. Laumon, L. Moret-Bailly -- \textit{Champs alg\'ebriques}, Springer-Verlag,
  (1999).}

\refitem{[Ly 1]}{S. Lysenko -- Local geometrized Rankin-Selberg method for GL($n$),
  \textit{Duke Math. J.} \textbf{111}, (2002), 451-493.}

\refitem{[Ly 2]}{S. Lysenko -- Global geometrized Rankin-Selberg method for GL($n$),
  http ://arxiv.org/abs/math.AG/0108208, (2001).}

\refitem{[Se]}{J.-P. Serre -- \textit{Groupes alg\'ebriques et corps de classes}, Hermann,
  (1975).}

\begin{flushright}
G\'erard Laumon\\
Universit\'e Paris-Sud\\
et CNRS, UMR 8628\\
Math\'ematique, B\^at. 425\\
F-91405 Orsay Cedex (France)\\
Gerard.Laumon@math.u-psud.fr
\end{flushright}


\end{document}
